\section*{الفصل \ref{ch:b2:himalaya-andes} --- مقارنة بين حزام تصادمي وحزام أنديزي}

\begin{solution}{exo:b2:himalaya-andes:1}
(أ) $r = h\rho_c/(\rho_m - \rho_c) = 3.0\times2700/600 = \qty{13.5}{km}$. (ب) $T = 35 + 3.0 +
13.5 = \qty{51.5}{km}$، أو $35 + 5.5\times3.0$. (ج) $h = \Delta T(\rho_m - \rho_c)/\rho_m =
10\times600/3300 = \qty{1.8}{km}$.
\end{solution}

\begin{solution}{exo:b2:himalaya-andes:2}
(أ) المنخفض الأمامي الهندي--الغانجي؛ الصدع الدفعي الجبهي الرئيسي؛ نطاق ما تحت الهيمالايا (سيواليك)؛ الصدع الدفعي
الحدّي الرئيسي؛ الهيمالايا الصغرى؛ الصدع الدفعي المركزي الرئيسي؛ متتالية الهيمالايا الكبرى؛ منظومة انفصال
جنوب التبت؛ متتالية الهيمالايا التيثيسية؛ ثم درز إندوس--تسانغبو والقشرة
الآسيوية. (ب) منظومة انفصال جنوب التبت. (ج) ترتفع الدرجة نحو الشمال: مولاس غير متحوّل،
فهيمالايا صغرى متوسطة الدرجة، فمتتالية هيمالايا كبرى عالية الدرجة ومنصهرة جزئيًا؛
وفوق STDS تهبط فجأة إلى رواسب تيثيسية غير متحوّلة أو ضعيفة
التحوّل، لأن الصدع العادي أزال الصخور التي بينهما.
\end{solution}

\begin{solution}{exo:b2:himalaya-andes:3}
(أ) $196.5\ \text{m}/9200\ \text{yr} = \qty{0.02136}{m}$ في السنة، أي نحو \qty{21.4}{mm} في السنة.
(ب) إنه أكثر بنسبة \qty{20}{\percent} من \qty{17.8}{mm} في السنة التي يعطيها GPS، لكنه من الرتبة نفسها:
فأحد المعدّلين متوسط على تسعة آلاف سنة، والآخر على عشرين.
(ج) يُستوعب التقصّر عبر السلسلة كله تقريبًا عند الجبهة: إذ إن MBT وMCT
وسائر الصدوع شمال MFT غير نشطة عمليًا عند خط الطول هذا؛ وكلها تتجذّر
في MHT نفسه، الذي يصعد إلى السطح عند MFT.
\end{solution}

\begin{solution}{exo:b2:himalaya-andes:4}
(أ) تصادمي: الدرز يدلّ على المحيط الزائل. (ب) من النمط الأنديزي: يحتاج القوس إلى كتلة مندسّة
نازلة. (ج) تصادمي في الغالب: مصهورات قشرية لقشرة متثخّنة مسخّنة (وتوجد
أيضًا في القشرة الأنديزية المتثخّنة، لكن صهارات القوس تسود هناك). (د) من النمط الأنديزي.
(ه) كلاهما: تلال سيواليك والمقدّمة تحت الأنديزية. (و) كلاهما: MHT في الهيمالايا،
وسطح الاندساس في الأنديز.
\end{solution}

\begin{solution}{exo:b2:himalaya-andes:5}
(أ) $17.8/45 = \qty{40}{\percent}$ في الشرق، و$20.5/45 = \qty{46}{\percent}$ في الغرب:
أي نحو النصف. (ب) إلى التثخّن وصدوع التبت وبقية آسيا، وإلى
بثق الكتل نحو الشرق على امتداد الصدوع المضربية الكبرى. (ج) $45 - 17.8 =
\qty{27.2}{mm}$ في السنة، أي \qty{272}{km} في \qty{10}{Myr}؛ و$45 - 20.5 = \qty{24.5}{mm}$ في
السنة، أي \qty{245}{km}.
\end{solution}

\begin{solution}{exo:b2:himalaya-andes:6}
(أ) $r = 4.5\times3.7 = \qty{16.65}{km}$؛ $\Delta F = 2700\times9.8\times3700\times(35\,000 +
1850 + 8325) = \qty{4.42e12}{N.m^{-1}}$. (ب) التبت: \qty{6.12e12}{N.m^{-1}}؛ وللألتيبلانو
\qty{72}{\percent} منه. (ج) مع $r = 4.5h$، يكون $\Delta F = \rho_c g\,(Hh + 2.75h^2)$: فمن أجل
$h \ll H$ يسود الحدّ $Hh$ ويكون $\Delta F$ متناسبًا مع $h$؛ ولهضبة عالية
ينمو الحدّ في $h^2$ أسرع. فكل كيلومتر إضافي من الارتفاع يكلّف دفعًا متزايدًا عند
الحواف، حتى لا تعود القوى المتاحة قادرة على رفع الهضبة؛ فتنبسط عندئذ
بدلًا من أن ترتفع.
\end{solution}

\begin{solution}{exo:b2:himalaya-andes:7}
(أ) $2700\times9.8\times8848.86/2 = \qty{1.17e8}{Pa} = \qty{117}{MPa}$. (ب) عند
\qty{100}{MPa}: $0.85\times100 = \qty{85}{MPa}$؛ وعند \qty{200}{MPa}: $0.85\times200 = 50 +
0.6\times200 = \qty{170}{MPa}$، فالفرعان يتّفقان. (ج) $h = 2\sigma_d/(\rho g) =
2\times1.70\times10^8/(2700\times9.8) = \qty{12.8}{km}$: رتبة مقدار، من رتبة
أعلى القمم، وليس أعلى منها بكثير.
\end{solution}

\begin{solution}{exo:b2:himalaya-andes:8}
(أ) $11\times111.19\times\cos\ang{30} = \qty{1059}{km}$. (ب) $9\times10^{-3}/1.059\times10^{6}
= \num{8.5e-9}$ في السنة (بين \num{6.6e-9} و\num{10.4e-9} من أجل 7 و\qty{11}{mm} في السنة)،
أي $8.5\times10^{-9}/3.156\times10^{7} = \qty{2.7e-16}{s^{-1}}$. (ج) المديان
متداخلان: $(12.4 \pm 4.7) \times 10^{-9}$ يمتدّ من 7.7 إلى \num{17.1e-9}. فجنوب الهضبة شمال
نيبال يتمدّد بسرعة الحزام ككلّ، أو أسرع منه قليلًا.
\end{solution}

\begin{solution}{exo:b2:himalaya-andes:9}
(أ) $0.9/6.6 = \qty{14}{\percent}$. (ب) $M_0 = 10^{1.5\times8.5 + 9.1} = 10^{21.85} =
\qty{7.1e21}{N.m}$. (ج) $7.1\times10^{21}/6.6\times10^{19} = 107$: زلزال واحد كهذا نحو
كل 110 سنوات تحت نيبال، وهو أكثر بكثير مما يسجّله التاريخ؛ فإمّا أن زلازل أكبر بكثير
مرتقبة، وإمّا أن جزءًا من العجز يتحرّر بطريقة أخرى.
\end{solution}

\begin{solution}{exo:b2:himalaya-andes:10}
(أ) خذ عمق التعويض عند قاعدة قشرة الألتيبلانو، $61.5 - 3.7 =
\qty{57.8}{km}$ تحت مستوى سطح البحر. وقشرة البونا تبلغ $50.5 - 4.4 = \qty{46.1}{km}$، فيحوي
عمود البونا \qty{11.7}{km} من الوشاح فوق ذلك العمق: $2800\times61.5 =
\rho_P\times50.5 + 3300\times11.7$. (ب) $\rho_P = (172\,200 - 38\,610)/50.5 =
\qty{2645}{kg.m^{-3}}$. (ج) الفرق، \qty{155}{kg.m^{-3}}، يقع عند الحدّ الأدنى من
القيم من 150 إلى \qty{200}{kg.m^{-3}} المذكورة: فالقشرة وحدها تتطلّب قشرة بونا خفيفة على نحو غير مألوف.
والعمود الأخفّ تحتها يحقّق ذلك بسهولة أكبر: غلاف صخري وشاحي أزاله التقشّر،
وحلّ محلّه أستينوسفير حارّ خفيف.
\end{solution}

\begin{solution}{exo:b2:himalaya-andes:11}
(أ) $g = 105/\tan\ang{30} = \qty{181.9}{km}$. (ب) $g = 90/\tan\ang{30} + 15/\tan\ang{4} =
155.9 + 214.5 = \qty{370.4}{km}$: ينتقل القوس \qty{188.5}{km} نحو الداخل. (ج) لا تبلغ الكتلة
عمق \qty{105}{km} تحت الصفيحة على مدى تلك \qty{300}{km}، ولا يقع فوقها إسفين وشاحي حارّ:
فلا تتكوّن صهارة ويموت القوس حتى تغوص الكتلة من جديد، أبعد نحو الشرق.
\end{solution}

\begin{solution}{exo:b2:himalaya-andes:12}
(أ) \qty{21}{mm} في السنة تساوي \qty{21}{km} في المليون سنة: $21\times37.5 =
\qty{787.5}{km^3}$ لكل كيلومتر من السلسلة في المليون سنة. (ب) $300/787.5 =
\qty{38}{\percent}$. (ج) في تثخّن السلسلة وجنوب التبت، وفي القشرة
المحمولة مع الهند إلى الأسفل تحت التبت، وفي المادة التي ينقلها MHT نحو الشمال.
\end{solution}

\begin{solution}{pb:b2:himalaya-andes:1}
\textbf{1.} سبع عيّنات من \ang{28.0} إلى \ang{37.0} شمالًا: $34\,550/7 = \qty{4936}{m}$، $h_T
= \qty{4.9}{km}$.
\textbf{2.} خمس عيّنات من \ang{-69} إلى \ang{-65}: $18\,650/5 = \qty{3730}{m}$، $h_A =
\qty{3.7}{km}$، وهي القيمة المذكورة للألتيبلانو؛ والبونا، الأبعد جنوبًا، أعلى.
\textbf{3.} التبت: على الأقل $9\times111.19 = \qty{1001}{km}$؛ والألتيبلانو: على الأقل $4\times
104.5 = \qty{418}{km}$.
\textbf{4.} $(5000 - 70)/(1.5\times111\,190) = 0.0296$: أي \qty{3.0}{\percent}.
\textbf{5.} يبدأ المقطع الأنديزي على قاع المحيط وفي الخندق، على كيلومترات تحت
مستوى سطح البحر، غرب الهضبة؛ ويبدأ المقطع التبتي على سهل قاري قرب مستوى سطح البحر،
وينتهي على حوض تاريم المرتفع.
\textbf{6.} $r = 4.5\times4.9 = \qty{22.05}{km}$؛ $T_T = 35 + 4.9 + 22.05 = \qty{62.0}{km}$.
\textbf{7.} $r = 4.5\times3.7 = \qty{16.65}{km}$؛ $T_A = 35 + 3.7 + 16.65 = \qty{55.4}{km}$.
\textbf{8.} التبت: \qty{62.0}{km} مقابل 70 إلى \qty{75}{km}، أي أقل بمقدار \qty{10.5}{km} من القيمة
الوسطى \qty{72.5}{km} (\qty{14}{\percent})؛ والألتيبلانو: \qty{55.4}{km} مقابل 59 إلى
\qty{64}{km}، أي أقل بمقدار \qty{6.1}{km} من \qty{61.5}{km} (\qty{10}{\percent}). فالتبت أبعد عن الواقع.
\textbf{9.} $72.5 - 35 = 4.9\times3300/(3300 - \rho_c)$: $3300 - \rho_c = 16\,170/37.5 =
431.2$، $\rho_c = \qty{2869}{kg.m^{-3}}$.
\textbf{10.} $26.5 = 3.7\times3300/(3300 - \rho_c)$: $3300 - \rho_c = 460.8$، $\rho_c =
\qty{2839}{kg.m^{-3}}$.
\textbf{11.} جزء من الحمل يحمله انثناء الصفيحة، لا جذر محلي؛ وربما لم يكن
سمك القشرة العادية \qty{35}{km}؛ وربما كان الوشاح تحتها أخفّ أو أكثف
مما افتُرض؛ وربما لا يكون الموهو الزلزالي قاعدة القشرة بمعنى الكثافة.
\textbf{12.} تقع قاعدة قشرة التبت على $72.5 - 4.9 = \qty{67.6}{km}$، وقاعدة قشرة الهند على
\qty{37.5}{km}: $t = \qty{30.1}{km}$. $\Delta g = 2\pi\times6.674\times10^{-11}\times370\times
30\,100 = \qty{4.67e-3}{m.s^{-2}}$، أي \qty{-467}{mGal} بالنسبة إلى الهند، قريبًا من
الفرق المرصود \qty{460}{mGal}.
\textbf{13.} $W_0 = 440\times60/35 = \qty{754.3}{km}$؛ $S = \qty{314.3}{km}$.
\textbf{14.} $(314.3 - 250)/314.3 = \qty{20}{\percent}$، وهو النقص الذي وُجد في الأنديز
الوسطى.
\textbf{15.} $0.55\times314.3 = \qty{173}{km}$، أي أقل من \qty{250}{km} التي تعطيها المقاطع
المتوازنة، التي تصير عندئذ أكثر من كافية.
\textbf{16.} يجب أن يقابل تقصّرَ الغطاء الرسوبي فوق \omterm{def:b2:alpine-orogen:decollement}{سطح الانفصال} تقصّرٌ في
القاعدة تحته، أبعد نحو الغرب، لأن كليهما يجب أن يفقد الطول نفسه عبر الحزام؛ وتقصّر
القاعدة يُثخّن القشرة تحت الكورديليرا الشرقية والهضبة.
\textbf{17.} $314.3/30 = \qty{10.5}{mm}$ في السنة، و$10.5/72 = \qty{15}{\percent}$ من
التقارب.
\textbf{18.} يستهلكه الاندساس: فصفيحة نازكا تغوص تحت القارة.
\textbf{19.} $r = \qty{22.05}{km}$؛ $\Delta F = 2700\times9.8\times4900\times(35\,000 + 2450
+ 11\,025) = \qty{6.28e12}{N.m^{-1}}$.
\textbf{20.} $2700\times9.8\times3700\times45\,175 = \qty{4.42e12}{N.m^{-1}}$؛ والنسبة
0.70.
\textbf{21.} $6.28\times10^{12}/10^{5} = \qty{6.3e7}{Pa}$، أي نحو \qty{63}{MPa}.
\textbf{22.} $9\times10^{-3}/1.06\times10^{6} = \num{8.5e-9}$ في السنة.
\textbf{23.} يعمل دفع الهند شمالًا--جنوبًا، وتوازنه في الهضبة العالية
وزنُ الهضبة الذاتي: فلا تستطيع الهضبة أن تتثخّن بعد ذلك شمالًا--جنوبًا، فيكون الإجهاد الرأسي،
أي الوزن، هو الأكبر؛ والإجهاد الأصغر شرقي--غربي، نحو الحواف الحرّة لآسيا،
فتتمدّد الهضبة في ذلك الاتجاه، على شقوق شمالية--جنوبية.
\textbf{24.} \textbf{\qty[text-series-to-math=true]{62}{km}}، أي نحو \qty{86}{\percent} من
\qty{72.5}{km} المقيسة.
\end{solution}
